\documentclass[12pt, a4paper]{article}

% ── Paquetes ────────────────────────────────────────────────
\usepackage[utf8]{inputenc}
\usepackage[T1]{fontenc}
\usepackage[spanish]{babel}
\usepackage{geometry}
\geometry{margin=2.5cm}
\usepackage{booktabs}
\usepackage{longtable}
\usepackage{array}
\usepackage{hyperref}
\usepackage{graphicx}
\usepackage{xcolor}
\usepackage{titlesec}
\usepackage{parskip}
\usepackage{enumitem}
\usepackage{fancyhdr}

% ── Encabezado / pie ────────────────────────────────────────
\pagestyle{fancy}
\fancyhf{}
\rhead{\small Bases de Datos · Sesión 1}
\lhead{\small UniGrades}
\cfoot{\thepage}

% ── Colores ─────────────────────────────────────────────────
\definecolor{unal}{RGB}{134, 31, 65}   % borgoña UNAL
\hypersetup{colorlinks=true, linkcolor=unal, urlcolor=unal}

% ── Título de secciones ─────────────────────────────────────
\titleformat{\section}{\large\bfseries\color{unal}}{}{0em}{}[\titlerule]
\titleformat{\subsection}{\normalsize\bfseries}{}{0em}{}

% ============================================================
\begin{document}

% ── PORTADA ─────────────────────────────────────────────────
\begin{titlepage}
    \centering
    \vspace*{2cm}
    {\LARGE\bfseries\color{unal} UniGrades\par}
    \vspace{0.5cm}
    {\large Gestor de Progreso Académico Universitario\par}
    \vspace{1.5cm}
    {\large\bfseries Proyecto Final — Bases de Datos\par}
    {\large Sesión 1: Kickoff y Diseño Conceptual\par}
    \vspace{2cm}

    \begin{tabular}{ll}
        \textbf{Integrante} & Sergio Chaves \\
        \textbf{Código}     & sechaves \\
    \end{tabular}

    \vfill
    {\large Universidad Nacional de Colombia\par}
    {\large Facultad de Ingeniería\par}
    \vspace{0.5cm}
    {\large \today\par}
\end{titlepage}

% ── 1. DESCRIPCIÓN DEL DOMINIO ──────────────────────────────
\section{Descripción del Dominio}

UniGrades es una aplicación web orientada a estudiantes universitarios que centraliza
el seguimiento de su progreso académico a lo largo de toda la carrera. El sistema
permite registrar las materias cursadas en cada período académico, configurar los
componentes de evaluación de cada asignatura (parciales, talleres, proyectos, etc.)
y registrar las calificaciones obtenidas, calculando automáticamente promedios por
materia, por semestre y el promedio global acumulado.

A diferencia de los sistemas institucionales (como el SIA de la UNAL), UniGrades
es una herramienta personal del estudiante: él mismo configura su plan de trabajo,
registra sus notas a medida que las recibe y visualiza su avance por tipología de
créditos. El sistema está diseñado para ser agnóstico a la institución, soportando
cualquier universidad, programa y escala de calificación.

Una base de datos relacional es indispensable porque el dominio involucra múltiples
entidades fuertemente relacionadas (universidad, programa, tipología, materia,
usuario, semestre, inscripción, componente y nota), datos históricos transaccionales
(semestres cursados, notas registradas por fecha), restricciones de integridad no
triviales (los porcentajes de un componente deben sumar 100, la nota mínima de
aprobación varía por materia) y consultas analíticas que cruzan varias dimensiones
(promedio ponderado por créditos, avance por tipología, roadmap sugerido vs.\ real).

% ── 2. REGLAS DE NEGOCIO ────────────────────────────────────
\section{Reglas de Negocio}

Las siguientes reglas describen las restricciones y comportamientos obligatorios del
sistema. Cada regla indica las entidades involucradas y la restricción que impone.

\begin{enumerate}[leftmargin=*, label=\textbf{RN-\arabic*.}]

    \item \textbf{Un estudiante pertenece a exactamente un programa.}\\
    \textit{Entidades:} \texttt{usuario}, \texttt{programa}.\\
    Un \texttt{usuario} debe tener asociado un único \texttt{programa} curricular al
    momento del registro. No puede existir un usuario sin programa asignado
    (\texttt{usuario\_programa\_id NOT NULL}).

    \item \textbf{Los porcentajes de los componentes de una materia deben sumar 100.}\\
    \textit{Entidades:} \texttt{componente}, \texttt{materia\_usuario}.\\
    La suma de \texttt{componente\_porcentaje} de todos los componentes asociados a
    una misma \texttt{materia\_usuario} debe ser exactamente 100. Esta restricción
    no puede expresarse con \texttt{NOT NULL} ni \texttt{UNIQUE}; requiere validación
    a nivel de aplicación o trigger.

    \item \textbf{Un estudiante no puede inscribir la misma materia dos veces en el mismo semestre.}\\
    \textit{Entidades:} \texttt{materia\_usuario}, \texttt{semestre}, \texttt{materia}.\\
    La combinación (\texttt{materia\_usuario\_usuario\_id},
    \texttt{materia\_usuario\_materia\_id},
    \texttt{materia\_usuario\_semestre\_id}) debe ser única. Un mismo estudiante
    puede cursar la misma materia en semestres distintos (habilitación/repetición),
    pero no dos veces en el mismo período.

    \item \textbf{El estado de una materia sigue transiciones definidas.}\\
    \textit{Entidades:} \texttt{materia\_usuario}.\\
    El atributo \texttt{materia\_usuario\_estado} solo puede tomar los valores
    \texttt{en\_curso}, \texttt{aprobada}, \texttt{reprobada} o \texttt{retirada}.
    Una materia con estado \texttt{aprobada} o \texttt{reprobada} no puede volver a
    \texttt{en\_curso} en el mismo registro; debe crearse una nueva inscripción en
    otro semestre.

    \item \textbf{El valor de una nota debe estar dentro del rango válido de la escala.}\\
    \textit{Entidades:} \texttt{nota}.\\
    \texttt{nota\_valor} debe estar entre 0.0 y 5.0 (escala colombiana estándar).
    Valores fuera de rango deben ser rechazados por la base de datos mediante una
    restricción \texttt{CHECK}.

    \item \textbf{Las tipologías de tipo nivelación no cuentan en el promedio acumulado.}\\
    \textit{Entidades:} \texttt{tipologia}, \texttt{materia}, \texttt{nota}.\\
    Las materias cuya tipología tiene \texttt{tipologia\_cuenta\_promedio = 0}
    (p.\ ej.\ Inglés I--IV en la UNAL) se excluyen del cálculo del promedio global
    y del promedio por semestre, aunque sí se registran y se validan.

    \item \textbf{Un componente acumula múltiples notas cuyo promedio define su calificación.}\\
    \textit{Entidades:} \texttt{componente}, \texttt{nota}.\\
    Un \texttt{componente} puede tener cero o más registros de \texttt{nota}
    (relación 1:N). La calificación del componente se calcula como el promedio
    aritmético de todas las notas registradas en ese momento. Esto permite modelar
    componentes acumulativos (p.\ ej.\ ``Talleres -- 20\%'') donde la cantidad
    de entregas no se conoce de antemano.

\end{enumerate}

% ── 3. DIAGRAMA E-R ─────────────────────────────────────────
\section{Diagrama Entidad-Relación}

El diagrama E-R fue elaborado en draw.io siguiendo la notación crow's foot.
Se incluye como archivo adjunto \texttt{Diagrama E-R.drawio} y su versión
exportada \texttt{Diagrama E-R.drawio.pdf} en el repositorio del proyecto.

\subsection*{Entidades del modelo}

El modelo cuenta con \textbf{9 entidades base}:

\begin{itemize}
    \item \textbf{UNIVERSIDAD} — maestra; almacena las instituciones soportadas.
    \item \textbf{PROGRAMA} — maestra; carrera curricular dentro de una universidad.
    \item \textbf{TIPOLOGIA} — maestra; categoría de materia dentro de un programa
          (Obligatoria, Optativa, Libre Elección, Nivelación, etc.).
    \item \textbf{MATERIA} — maestra; asignatura del programa con semestre sugerido (roadmap).
    \item \textbf{USUARIO} — maestra; el estudiante registrado en la plataforma.
    \item \textbf{SEMESTRE} — transaccional; período académico real cursado por el estudiante.
    \item \textbf{MATERIA\_USUARIO} — asociativa ternaria (M:N con atributos);
          vincula usuario, materia y semestre. Registra el estado de la materia.
    \item \textbf{COMPONENTE} — transaccional; evaluación configurada por el estudiante
          (nombre, porcentaje, nota mínima).
    \item \textbf{NOTA} — transaccional; calificación individual dentro de un componente.
\end{itemize}

\subsection*{Relaciones M:N con atributos propios}

\begin{enumerate}
    \item \textbf{MATERIA\_USUARIO} (USUARIO × MATERIA × SEMESTRE): tabla asociativa
    ternaria. Atributos propios: \texttt{materia\_usuario\_estado}
    (\texttt{en\_curso / aprobada / reprobada / retirada}).
    \item \textbf{COMPONENTE} como extensión de MATERIA\_USUARIO: cada inscripción
    genera sus propios componentes de evaluación con atributos
    \texttt{componente\_porcentaje} y \texttt{componente\_nota\_minima}.
\end{enumerate}

\subsection*{Jerarquía / auto-relación}

La entidad \textbf{TIPOLOGIA} tiene el atributo \texttt{tipologia\_cuenta\_promedio}
que establece una jerarquía funcional entre tipologías: aquellas con valor 0 son de
nivelación y se comportan de forma distinta en los cálculos. Adicionalmente,
\textbf{MATERIA} referencia a \textbf{PROGRAMA} y a \textbf{TIPOLOGIA} de forma
simultánea, creando una dependencia cruzada que modela la estructura curricular
jerárquica (universidad → programa → tipología → materia).

\subsection*{Datos temporales}

La entidad \textbf{SEMESTRE} registra \texttt{semestre\_year} y
\texttt{semestre\_periodo}, permitiendo consultas históricas por rango de períodos
académicos. La entidad \textbf{NOTA} registra \texttt{nota\_fecha\_registro},
habilitando consultas cronológicas del progreso del estudiante.

% ── 4. DICCIONARIO DE DATOS ─────────────────────────────────
\section{Diccionario de Datos Preliminar}

\begingroup
\small
\setlength{\tabcolsep}{4pt}

% ── UNIVERSIDAD ─────────────────────────────────────────────
\subsection*{UNIVERSIDAD}
\begin{longtable}{>{\ttfamily}p{4.5cm} p{2.5cm} p{3.2cm} p{4.5cm}}
\toprule
\normalfont\textbf{Atributo} & \textbf{Tipo} & \textbf{Restricción} & \textbf{Descripción} \\
\midrule
\endfirsthead
\toprule
\normalfont\textbf{Atributo} & \textbf{Tipo} & \textbf{Restricción} & \textbf{Descripción} \\
\midrule
\endhead
universidad\_id         & INT UNSIGNED   & PK, NOT NULL, AUTO   & Identificador único de la universidad. \\
universidad\_nombre     & VARCHAR(150)   & NOT NULL             & Nombre oficial de la institución. \\
universidad\_sigla      & VARCHAR(20)    & NULL                 & Sigla o acrónimo (ej. UNAL). \\
universidad\_pais       & VARCHAR(80)    & NOT NULL, DEF Colombia & País sede de la institución. \\
universidad\_ciudad     & VARCHAR(80)    & NULL                 & Ciudad sede. \\
universidad\_logo\_url  & VARCHAR(255)   & NULL                 & URL del logotipo para la interfaz. \\
universidad\_created\_at & TIMESTAMP     & DEF NOW()            & Fecha y hora de registro. \\
\bottomrule
\end{longtable}

% ── PROGRAMA ────────────────────────────────────────────────
\subsection*{PROGRAMA}
\begin{longtable}{>{\ttfamily}p{4.5cm} p{2.5cm} p{3.2cm} p{4.5cm}}
\toprule
\normalfont\textbf{Atributo} & \textbf{Tipo} & \textbf{Restricción} & \textbf{Descripción} \\
\midrule
\endfirsthead
\toprule
\normalfont\textbf{Atributo} & \textbf{Tipo} & \textbf{Restricción} & \textbf{Descripción} \\
\midrule
\endhead
programa\_id                & INT UNSIGNED   & PK, NOT NULL, AUTO   & Identificador único del programa. \\
programa\_universidad\_id   & INT UNSIGNED   & FK, NOT NULL         & Universidad a la que pertenece. \\
programa\_nombre            & VARCHAR(150)   & NOT NULL             & Nombre del programa curricular. \\
programa\_facultad          & VARCHAR(150)   & NULL                 & Facultad que lo administra. \\
programa\_total\_creditos   & SMALLINT       & NOT NULL, DEF 0      & Total de créditos para graduarse. \\
programa\_created\_at       & TIMESTAMP      & DEF NOW()            & Fecha de registro. \\
\bottomrule
\end{longtable}

% ── TIPOLOGIA ───────────────────────────────────────────────
\subsection*{TIPOLOGIA}
\begin{longtable}{>{\ttfamily}p{4.5cm} p{2.5cm} p{3.2cm} p{4.5cm}}
\toprule
\normalfont\textbf{Atributo} & \textbf{Tipo} & \textbf{Restricción} & \textbf{Descripción} \\
\midrule
\endfirsthead
\toprule
\normalfont\textbf{Atributo} & \textbf{Tipo} & \textbf{Restricción} & \textbf{Descripción} \\
\midrule
\endhead
tipologia\_id                   & INT UNSIGNED  & PK, NOT NULL, AUTO   & Identificador único. \\
tipologia\_programa\_id         & INT UNSIGNED  & FK, NOT NULL         & Programa al que pertenece. \\
tipologia\_nombre               & VARCHAR(100)  & NOT NULL             & Nombre de la tipología (ej. Disciplinar Obligatoria). \\
tipologia\_creditos\_requeridos & SMALLINT      & NOT NULL, DEF 0      & Créditos mínimos que debe completar el estudiante en esta tipología. \\
tipologia\_cuenta\_promedio     & TINYINT(1)    & NOT NULL, DEF 1      & 1 = cuenta en el promedio; 0 = no cuenta (ej. Nivelación). \\
tipologia\_created\_at          & TIMESTAMP     & DEF NOW()            & Fecha de registro. \\
\bottomrule
\end{longtable}

% ── MATERIA ─────────────────────────────────────────────────
\subsection*{MATERIA}
\begin{longtable}{>{\ttfamily}p{4.5cm} p{2.5cm} p{3.2cm} p{4.5cm}}
\toprule
\normalfont\textbf{Atributo} & \textbf{Tipo} & \textbf{Restricción} & \textbf{Descripción} \\
\midrule
\endfirsthead
\toprule
\normalfont\textbf{Atributo} & \textbf{Tipo} & \textbf{Restricción} & \textbf{Descripción} \\
\midrule
\endhead
materia\_id                         & INT UNSIGNED  & PK, NOT NULL, AUTO        & Identificador único. \\
materia\_programa\_id               & INT UNSIGNED  & FK, NOT NULL              & Programa al que pertenece la materia. \\
materia\_tipologia\_id              & INT UNSIGNED  & FK, NOT NULL              & Tipología de la materia. \\
materia\_codigo                     & VARCHAR(30)   & NULL, UNIQUE por programa & Código libre según la universidad (ej. 1000003-B). \\
materia\_nombre                     & VARCHAR(150)  & NOT NULL                  & Nombre oficial de la asignatura. \\
materia\_creditos                   & TINYINT       & NOT NULL, DEF 3           & Número de créditos académicos. \\
materia\_nota\_minima\_aprobacion   & DECIMAL(3,1)  & NOT NULL, DEF 3.0         & Nota mínima para aprobar la materia. \\
materia\_semestre\_sugerido         & TINYINT       & NULL                      & Semestre recomendado en el roadmap del programa. \\
materia\_created\_at                & TIMESTAMP     & DEF NOW()                 & Fecha de registro. \\
\bottomrule
\end{longtable}

% ── USUARIO ─────────────────────────────────────────────────
\subsection*{USUARIO}
\begin{longtable}{>{\ttfamily}p{4.5cm} p{2.5cm} p{3.2cm} p{4.5cm}}
\toprule
\normalfont\textbf{Atributo} & \textbf{Tipo} & \textbf{Restricción} & \textbf{Descripción} \\
\midrule
\endfirsthead
\toprule
\normalfont\textbf{Atributo} & \textbf{Tipo} & \textbf{Restricción} & \textbf{Descripción} \\
\midrule
\endhead
usuario\_id                  & INT UNSIGNED  & PK, NOT NULL, AUTO   & Identificador único del estudiante. \\
usuario\_programa\_id        & INT UNSIGNED  & FK, NOT NULL         & Programa en el que está matriculado. \\
usuario\_nombre              & VARCHAR(100)  & NOT NULL             & Nombre(s) del estudiante. \\
usuario\_apellido            & VARCHAR(100)  & NOT NULL             & Apellido(s) del estudiante. \\
usuario\_email               & VARCHAR(150)  & NOT NULL, UNIQUE     & Correo electrónico de acceso. \\
usuario\_password\_hash      & VARCHAR(255)  & NOT NULL             & Contraseña cifrada (hash). \\
usuario\_codigo\_estudiantil & VARCHAR(20)   & NULL                 & Código estudiantil institucional. \\
usuario\_avatar\_url         & VARCHAR(255)  & NULL                 & URL de foto de perfil. \\
usuario\_created\_at         & TIMESTAMP     & DEF NOW()            & Fecha de registro en la plataforma. \\
\bottomrule
\end{longtable}

% ── SEMESTRE ────────────────────────────────────────────────
\subsection*{SEMESTRE}
\begin{longtable}{>{\ttfamily}p{4.5cm} p{2.5cm} p{3.2cm} p{4.5cm}}
\toprule
\normalfont\textbf{Atributo} & \textbf{Tipo} & \textbf{Restricción} & \textbf{Descripción} \\
\midrule
\endfirsthead
\toprule
\normalfont\textbf{Atributo} & \textbf{Tipo} & \textbf{Restricción} & \textbf{Descripción} \\
\midrule
\endhead
semestre\_id         & INT UNSIGNED   & PK, NOT NULL, AUTO              & Identificador único. \\
semestre\_usuario\_id & INT UNSIGNED  & FK, NOT NULL                    & Estudiante al que pertenece el semestre. \\
semestre\_numero     & TINYINT        & NOT NULL                        & Número de semestre cursado por el estudiante (1, 2, 3\ldots). \\
semestre\_year       & YEAR           & NOT NULL                        & Año académico (ej. 2025). \\
semestre\_periodo    & TINYINT        & NOT NULL, CHECK (1 OR 2)        & Período: 1 = primer semestre, 2 = segundo semestre. \\
semestre\_created\_at & TIMESTAMP     & DEF NOW()                       & Fecha de creación del registro. \\
\bottomrule
\end{longtable}

% ── MATERIA_USUARIO ─────────────────────────────────────────
\subsection*{MATERIA\_USUARIO}
\begin{longtable}{>{\ttfamily}p{4.8cm} p{2.5cm} p{3.2cm} p{4.2cm}}
\toprule
\normalfont\textbf{Atributo} & \textbf{Tipo} & \textbf{Restricción} & \textbf{Descripción} \\
\midrule
\endfirsthead
\toprule
\normalfont\textbf{Atributo} & \textbf{Tipo} & \textbf{Restricción} & \textbf{Descripción} \\
\midrule
\endhead
materia\_usuario\_id           & INT UNSIGNED & PK, NOT NULL, AUTO           & Identificador único de la inscripción. \\
materia\_usuario\_usuario\_id  & INT UNSIGNED & FK, NOT NULL                 & Estudiante que cursa la materia. \\
materia\_usuario\_materia\_id  & INT UNSIGNED & FK, NOT NULL                 & Materia inscrita. \\
materia\_usuario\_semestre\_id & INT UNSIGNED & FK, NOT NULL                 & Semestre en que se cursa. \\
materia\_usuario\_estado       & ENUM         & NOT NULL, DEF en\_curso      & Estado: en\_curso / aprobada / reprobada / retirada. \\
materia\_usuario\_created\_at  & TIMESTAMP    & DEF NOW()                    & Fecha de registro de la inscripción. \\
\bottomrule
\end{longtable}

% ── COMPONENTE ──────────────────────────────────────────────
\subsection*{COMPONENTE}
\begin{longtable}{>{\ttfamily}p{4.5cm} p{2.5cm} p{3.2cm} p{4.5cm}}
\toprule
\normalfont\textbf{Atributo} & \textbf{Tipo} & \textbf{Restricción} & \textbf{Descripción} \\
\midrule
\endfirsthead
\toprule
\normalfont\textbf{Atributo} & \textbf{Tipo} & \textbf{Restricción} & \textbf{Descripción} \\
\midrule
\endhead
componente\_id                   & INT UNSIGNED  & PK, NOT NULL, AUTO   & Identificador único del componente. \\
componente\_materia\_usuario\_id & INT UNSIGNED  & FK, NOT NULL         & Inscripción a la que pertenece el componente. \\
componente\_nombre               & VARCHAR(100)  & NOT NULL             & Nombre del componente (ej. Parcial 01, Talleres). \\
componente\_porcentaje           & DECIMAL(5,2)  & NOT NULL             & Peso del componente en la nota final (0.00--100.00). \\
componente\_nota\_minima         & DECIMAL(3,1)  & NULL                 & Nota mínima para superar el componente (opcional). \\
componente\_orden                & TINYINT       & NOT NULL, DEF 1      & Orden de visualización del componente. \\
componente\_created\_at          & TIMESTAMP     & DEF NOW()            & Fecha de creación. \\
\bottomrule
\end{longtable}

% ── NOTA ────────────────────────────────────────────────────
\subsection*{NOTA}
\begin{longtable}{>{\ttfamily}p{4.5cm} p{2.5cm} p{3.2cm} p{4.5cm}}
\toprule
\normalfont\textbf{Atributo} & \textbf{Tipo} & \textbf{Restricción} & \textbf{Descripción} \\
\midrule
\endfirsthead
\toprule
\normalfont\textbf{Atributo} & \textbf{Tipo} & \textbf{Restricción} & \textbf{Descripción} \\
\midrule
\endhead
nota\_id              & INT UNSIGNED  & PK, NOT NULL, AUTO             & Identificador único de la nota. \\
nota\_componente\_id  & INT UNSIGNED  & FK, NOT NULL                   & Componente al que pertenece la nota. \\
nota\_nombre          & VARCHAR(100)  & NOT NULL                       & Identificador de la entrega (ej. Taller 1, Quiz 3). \\
nota\_valor           & DECIMAL(3,1)  & NOT NULL, CHECK (0.0--5.0)     & Calificación obtenida en escala 0.0--5.0. \\
nota\_fecha\_registro & DATE          & NULL                           & Fecha en que se recibió o registró la nota. \\
nota\_observacion     & VARCHAR(255)  & NULL                           & Comentario opcional del estudiante. \\
nota\_created\_at     & TIMESTAMP     & DEF NOW()                      & Fecha de creación del registro. \\
\bottomrule
\end{longtable}

\endgroup

\end{document}
